% Основной текст отчёта.
% Графики из code/lab1.ipynb сохраняйте в папку tex/images
% под именами, указанными в плейсхолдерах ниже.

\newcommand{\imageplaceholder}[3]{%
  \begin{figure}[htbp]
    \centering
    \IfFileExists{images/#1}{%
      \includegraphics[width=0.82\textwidth]{#1}%
    }{%
      \fbox{%
        \parbox[c][5cm][c]{0.78\textwidth}{%
          \centering
          Место для рисунка\\[0.4em]
          \texttt{images/\detokenize{#1}}
        }%
      }%
    }
    \caption{#2}
    \label{#3}
  \end{figure}
}

\tableofcontents
\newpage

\section{Цель работы}

Исследовать положения равновесия нелинейных динамических систем, выполнить
их линеаризацию, определить тип и устойчивость изолированных положений
равновесия, построить фазовые портреты и синтезировать локально
стабилизирующие регуляторы.

\section{Исследование автономных нелинейных систем}

\subsection{Система 1}

\begin{equation}
  \begin{cases}
    \dot x_1=-x_1+2x_1^3+x_2,\\
    \dot x_2=-x_1-x_2.
  \end{cases}
  \label{eq:system-1}
\end{equation}

Точки равновесия находятся из уравнений
\begin{equation}
  \begin{cases}
    -x_1+2x_1^3+x_2=0,\\
    -x_1-x_2=0.
  \end{cases}
\end{equation}
Матрица Якоби имеет вид
\begin{equation}
  A_1(x_1,x_2)=
  \begin{pmatrix}
    -1+6x_1^2 & 1\\
    -1 & -1
  \end{pmatrix}.
\end{equation}

\imageplaceholder{phase_system_1.png}{Фазовый портрет системы 1}{fig:phase-system-1}

\subsection{Система 2}

\begin{equation}
  \begin{cases}
    \dot x_1=x_1+x_1x_2,\\
    \dot x_2=-x_2+x_2^2+x_1x_2-x_1^3.
  \end{cases}
  \label{eq:system-2}
\end{equation}

Точки равновесия находятся из уравнений
\begin{equation}
  \begin{cases}
    x_1+x_1x_2=0,\\
    -x_2+x_2^2+x_1x_2-x_1^3=0.
  \end{cases}
\end{equation}
Матрица Якоби:
\begin{equation}
  A_2(x_1,x_2)=
  \begin{pmatrix}
    1+x_2 & x_1\\
    x_2-3x_1^2 & -1+2x_2+x_1
  \end{pmatrix}.
\end{equation}

\imageplaceholder{phase_system_2.png}{Фазовый портрет системы 2}{fig:phase-system-2}

\subsection{Система 3}

\begin{equation}
  \begin{cases}
    \dot x_1=x_2,\\
    \dot x_2=-x_1+x_2\left(1-x_1^2+0{,}1x_1^4\right).
  \end{cases}
  \label{eq:system-3}
\end{equation}

Точки равновесия находятся из уравнений
\begin{equation}
  \begin{cases}
    x_2=0,\\
    -x_1+x_2\left(1-x_1^2+0{,}1x_1^4\right)=0.
  \end{cases}
\end{equation}
Матрица Якоби:
\begin{equation}
  A_3(x_1,x_2)=
  \begin{pmatrix}
    0 & 1\\
    -1+x_2(-2x_1+0{,}4x_1^3) & 1-x_1^2+0{,}1x_1^4
  \end{pmatrix}.
\end{equation}

\imageplaceholder{phase_system_3.png}{Фазовый портрет системы 3}{fig:phase-system-3}

\subsection{Система 4}

\begin{equation}
  \begin{cases}
    \dot x_1=(x_1-x_2)(1-x_1^2-x_2^2),\\
    \dot x_2=(x_1+x_2)(1-x_1^2-x_2^2).
  \end{cases}
  \label{eq:system-4}
\end{equation}

Точки равновесия находятся из уравнений
\begin{equation}
  \begin{cases}
    (x_1-x_2)(1-x_1^2-x_2^2)=0,\\
    (x_1+x_2)(1-x_1^2-x_2^2)=0.
  \end{cases}
\end{equation}
Обозначив $q=1-x_1^2-x_2^2$, получим матрицу Якоби
\begin{equation}
  A_4(x_1,x_2)=
  \begin{pmatrix}
    q-2x_1(x_1-x_2) & -q-2x_2(x_1-x_2)\\
    q-2x_1(x_1+x_2) & q-2x_2(x_1+x_2)
  \end{pmatrix}.
\end{equation}

Для дополнительного исследования используется переход
\begin{equation}
  x_1=r\cos\varphi,\qquad x_2=r\sin\varphi,
\end{equation}
где
\begin{equation}
  \dot r=\frac{x_1\dot x_1+x_2\dot x_2}{r},\qquad
  \dot\varphi=\frac{x_1\dot x_2-x_2\dot x_1}{r^2}.
\end{equation}
После подстановки системы~\eqref{eq:system-4} необходимо исследовать уравнения
\begin{equation}
  \dot r=r(1-r^2),\qquad \dot\varphi=1-r^2.
\end{equation}

\textbf{Вывод об устойчивости при различных значениях $r$:} \ldots

\imageplaceholder{phase_system_4.png}{Фазовый портрет системы 4}{fig:phase-system-4}

\subsection{Система 5}

\begin{equation}
  \begin{cases}
    \dot x_1=-x_1^3+x_2,\\
    \dot x_2=x_1-x_2^3.
  \end{cases}
  \label{eq:system-5}
\end{equation}

Точки равновесия находятся из уравнений
\begin{equation}
  \begin{cases}
    -x_1^3+x_2=0,\\
    x_1-x_2^3=0.
  \end{cases}
\end{equation}
Матрица Якоби:
\begin{equation}
  A_5(x_1,x_2)=
  \begin{pmatrix}
    -3x_1^2 & 1\\
    1 & -3x_2^2
  \end{pmatrix}.
\end{equation}

\imageplaceholder{phase_system_5.png}{Фазовый портрет системы 5}{fig:phase-system-5}

\subsection{Система 6}

\begin{equation}
  \begin{cases}
    \dot x_1=-x_1^3+x_2^3,\\
    \dot x_2=x_2^3x_1-x_2^3.
  \end{cases}
  \label{eq:system-6}
\end{equation}

Точки равновесия находятся из уравнений
\begin{equation}
  \begin{cases}
    -x_1^3+x_2^3=0,\\
    x_2^3(x_1-1)=0.
  \end{cases}
\end{equation}
Матрица Якоби:
\begin{equation}
  A_6(x_1,x_2)=
  \begin{pmatrix}
    -3x_1^2 & 3x_2^2\\
    x_2^3 & 3x_2^2(x_1-1)
  \end{pmatrix}.
\end{equation}

\imageplaceholder{phase_system_6.png}{Фазовый портрет системы 6}{fig:phase-system-6}

\subsection{Система 7}

\begin{equation}
  \begin{cases}
    \dot x_1=-x_1^3+x_2^3,\\
    \dot x_2=x_1+3x_3-x_2^3,\\
    \dot x_3=x_1x_3-x_2^3-\sin x_1.
  \end{cases}
  \label{eq:system-7}
\end{equation}

Точки равновесия находятся из уравнений
\begin{equation}
  \begin{cases}
    -x_1^3+x_2^3=0,\\
    x_1+3x_3-x_2^3=0,\\
    x_1x_3-x_2^3-\sin x_1=0.
  \end{cases}
\end{equation}
Матрица Якоби:
\begin{equation}
  A_7(x_1,x_2,x_3)=
  \begin{pmatrix}
    -3x_1^2 & 3x_2^2 & 0\\
    1 & -3x_2^2 & 3\\
    x_3-\cos x_1 & -3x_2^2 & x_1
  \end{pmatrix}.
\end{equation}

Для системы 7 построение фазового портрета заданием не требуется.

\section{Синтез локально стабилизирующих регуляторов}

Для управляемой системы $\dot{\bm x}=\bm f(\bm x,\bm u)$ необходимо:
\begin{enumerate}[label=\arabic*)]
  \item найти пары $\left(\bm x_{\mathrm{eq}},\bm u_{\mathrm{eq}}\right)$,
  удовлетворяющие $\bm f(\bm x_{\mathrm{eq}},\bm u_{\mathrm{eq}})=0$;
  \item построить линейную модель
  \begin{equation}
    \Delta\dot{\bm x}=A\Delta\bm x+B\Delta\bm u,
    \qquad
    A=\left.\frac{\partial\bm f}{\partial\bm x}\right|_{\mathrm{eq}},
    \quad
    B=\left.\frac{\partial\bm f}{\partial\bm u}\right|_{\mathrm{eq}};
  \end{equation}
  \item проверить управляемость пары $(A,B)$;
  \item выбрать желаемые полюса и найти матрицу обратной связи $K$;
  \item использовать закон управления
  \begin{equation}
    \bm u=\bm u_{\mathrm{eq}}-K(\bm x-\bm x_{\mathrm{eq}})
  \end{equation}
  и проверить устойчивость замкнутой системы численно.
\end{enumerate}

\subsection{Управляемая система 1}

\begin{equation}
  \begin{cases}
    \dot x_1=-x_1+2x_1^3+x_2+\sin u_1,\\
    \dot x_2=-x_1-x_2+3\sin u_2.
  \end{cases}
  \label{eq:controlled-system-1}
\end{equation}

Условия равновесия:
\begin{equation}
  \begin{cases}
    -x_1+2x_1^3+x_2+\sin u_1=0,\\
    -x_1-x_2+3\sin u_2=0.
  \end{cases}
\end{equation}
Матрицы линеаризованной системы:
\begin{equation}
  A=
  \begin{pmatrix}
    -1+6x_1^2 & 1\\
    -1 & -1
  \end{pmatrix},
  \qquad
  B=
  \begin{pmatrix}
    \cos u_1 & 0\\
    0 & 3\cos u_2
  \end{pmatrix}.
\end{equation}

\textbf{Выбранная точка равновесия:}
$\bm x_{\mathrm{eq}}=\ldots$, $\bm u_{\mathrm{eq}}=\ldots$.

\textbf{Желаемые полюса:} $\ldots$.

\textbf{Матрица регулятора:} $K=\ldots$.

\imageplaceholder{control_system_1.png}{Переходные процессы замкнутой управляемой системы 1}{fig:control-system-1}

\subsection{Управляемая система 2}

\begin{equation}
  \begin{cases}
    \dot x_1=x_2+x_1x_2+u^3,\\
    \dot x_2=-x_2+x_2^2-x_1^3+\sin u.
  \end{cases}
  \label{eq:controlled-system-2}
\end{equation}

Условия равновесия:
\begin{equation}
  \begin{cases}
    x_2+x_1x_2+u^3=0,\\
    -x_2+x_2^2-x_1^3+\sin u=0.
  \end{cases}
\end{equation}
Матрицы линеаризованной системы:
\begin{equation}
  A=
  \begin{pmatrix}
    x_2 & 1+x_1\\
    -3x_1^2 & -1+2x_2
  \end{pmatrix},
  \qquad
  B=
  \begin{pmatrix}
    3u^2\\
    \cos u
  \end{pmatrix}.
\end{equation}

\textbf{Выбранная точка равновесия:}
$\bm x_{\mathrm{eq}}=\ldots$, $u_{\mathrm{eq}}=\ldots$.

\textbf{Желаемые полюса:} $\ldots$.

\textbf{Коэффициенты регулятора:} $K=\ldots$.

\imageplaceholder{control_system_2.png}{Переходные процессы замкнутой управляемой системы 2}{fig:control-system-2}


\section{Вывод}

В данном разделе необходимо сопоставить аналитическую классификацию точек
равновесия с фазовыми портретами, сформулировать вывод об устойчивости
системы 4 и оценить качество работы синтезированных регуляторов.
